\documentclass[11pt,a4paper]{article}

% Packages
\usepackage[margin=2.5cm]{geometry}
\usepackage{microtype}
\usepackage{graphicx}
\usepackage{tikz}
\usetikzlibrary{shapes,arrows.meta,positioning,calc,fit,backgrounds,shadows.blur,decorations.pathmorphing,decorations.pathreplacing}
\usepackage{colortbl}
\usepackage{enumitem}
\usepackage{booktabs}
\usepackage{longtable}
\usepackage{hyperref}
\usepackage{xcolor}
\usepackage{fancyhdr}
\usepackage{titlesec}
\usepackage[most]{tcolorbox}
\usepackage{multicol}
\usepackage{amsmath,amssymb}
\usepackage{multirow}
\usepackage{setspace}
\usepackage{fontawesome5}
\usepackage{listings}

% Spacing
\setlength{\parskip}{3pt plus 1pt minus 1pt}
\setlength{\parindent}{0pt}

% Colors
\definecolor{primary}{HTML}{0E4D64}
\definecolor{secondary}{HTML}{137177}
\definecolor{accent}{HTML}{C84B31}
\definecolor{lightbg}{HTML}{E8F6F3}
\definecolor{defbg}{HTML}{FDF2E9}
\definecolor{greenhl}{HTML}{2E8B57}
\definecolor{codebg}{HTML}{F4F6F7}

% Hyperref
\hypersetup{colorlinks=true, linkcolor=primary, urlcolor=secondary, citecolor=primary}

% Listing style
\lstset{
  backgroundcolor=\color{codebg},
  basicstyle=\ttfamily\small,
  keywordstyle=\color{primary}\bfseries,
  stringstyle=\color{greenhl},
  commentstyle=\color{gray},
  breaklines=true,
  frame=l,
  framerule=2pt,
  rulecolor=\color{secondary!40},
  xleftmargin=8pt,
  framexleftmargin=6pt,
  aboveskip=6pt,
  belowskip=6pt
}

\newcolumntype{L}[1]{>{\raggedright\arraybackslash}p{#1}}
\newcolumntype{C}[1]{>{\centering\arraybackslash}p{#1}}

% Header/Footer
\pagestyle{fancy}
\fancyhf{}
\fancyhead[L]{\small\textcolor{primary}{\textbf{Synopsis Guide and Demo --- MSIT}}}
\fancyhead[R]{\small\textcolor{primary}{Lecture Handout}}
\fancyfoot[C]{\small\thepage}
\renewcommand{\headrulewidth}{0.5pt}
\renewcommand{\footrulewidth}{0.3pt}

% Section formatting
\titleformat{\section}{\Large\bfseries\color{primary}}{\thesection}{0.6em}{}[\vspace{2pt}\titlerule]
\titleformat{\subsection}{\large\bfseries\color{secondary}}{\thesubsection}{0.5em}{}
\titleformat{\subsubsection}{\normalsize\bfseries\color{primary}}{\thesubsubsection}{0.5em}{}
\titlespacing*{\section}{0pt}{14pt}{8pt}
\titlespacing*{\subsection}{0pt}{10pt}{5pt}

% Custom boxes
\newtcolorbox{definitionbox}[1][] {
  enhanced,
  colback=defbg, colframe=orange!60!black,
  colbacktitle=orange!70!black, coltitle=white,
  title={\faIcon{book}~#1}, fonttitle=\bfseries,
  boxrule=0pt, leftrule=4pt, arc=2pt,
  shadow={1mm}{-1mm}{0mm}{black!15},
  top=4pt, bottom=4pt, left=8pt, right=8pt
}
\newtcolorbox{conceptbox}[1][] {
  enhanced,
  colback=lightbg, colframe=secondary,
  colbacktitle=secondary, coltitle=white,
  title={\faIcon{lightbulb}~#1}, fonttitle=\bfseries,
  boxrule=0pt, leftrule=4pt, arc=2pt,
  shadow={1mm}{-1mm}{0mm}{black!15},
  top=4pt, bottom=4pt, left=8pt, right=8pt
}
\newtcolorbox{alertbox}[1][] {
  enhanced,
  colback=red!4, colframe=accent,
  colbacktitle=accent, coltitle=white,
  title={\faIcon{exclamation-triangle}~#1}, fonttitle=\bfseries,
  boxrule=0pt, leftrule=4pt, arc=2pt,
  shadow={1mm}{-1mm}{0mm}{black!12},
  top=3pt, bottom=3pt, left=8pt, right=8pt
}
\newtcolorbox{examplebox}[1][] {
  enhanced,
  colback=green!4, colframe=greenhl,
  colbacktitle=greenhl, coltitle=white,
  title={\faIcon{code}~#1}, fonttitle=\bfseries,
  boxrule=0pt, leftrule=4pt, arc=2pt,
  shadow={1mm}{-1mm}{0mm}{black!12},
  top=3pt, bottom=3pt, left=8pt, right=8pt
}

\begin{document}

% Title
\begin{center}
\vspace*{0.5cm}
{\Huge\bfseries\textcolor{primary}{Synopsis Guide and Demo Synopsis}}\\[8pt]
{\Large\bfseries\textcolor{primary}{Topics in IT --- MSIT}}\\[12pt]
{\large\textcolor{secondary}{Based on IUB Guidelines for Proposal/Synopses and Theses Preparation 2023}}\\[14pt]
{\color{secondary}\rule{0.78\textwidth}{1.5pt}}\\[12pt]
{\Large\textsc{MSIT Lecture Handout}}\\[6pt]
{\large\textcolor{gray!70!black}{Course: Topics in IT \quad Program: MSIT}}\\[8pt]
{\normalsize\textcolor{gray!70!black}{Synopsis Organization \textbullet\ Section-by-Section Guide \textbullet\ Demo Synopsis \textbullet\ APA References}}
\vspace{0.3cm}
\end{center}

\tableofcontents
\newpage

%=============================================================
\section{Purpose of This Handout}
%=============================================================

\begin{definitionbox}[What This Handout Is]
This handout is a practical guide for MSIT students preparing a synopsis in the course \emph{Topics in IT}. It follows the organizational pattern visible in the IUB synopsis summary chart and turns that pattern into a usable classroom guide plus an illustrative demo synopsis.
\end{definitionbox}

\begin{conceptbox}[How to Use It]
Read this handout in two ways. First, use the guide sections to understand what each part of a synopsis should contain. Second, use the demo synopsis as a model for tone, sequence, heading structure, and section depth.
\end{conceptbox}

\begin{alertbox}[Important Note]
This handout is \textbf{based on IUB Guidelines for Proposal/Synopses and Theses Preparation 2023}. It is a classroom adaptation for MSIT students, not a replacement for the official university document. When preparing your final submission, always cross-check departmental and supervisor requirements.
\end{alertbox}

%=============================================================
\section{Summary Chart for Synopsis Organization}
%=============================================================

\begin{center}
{\Large\bfseries\textcolor{primary}{Summary Chart for the Organization of the Synopsis and its Style and Form}}
\end{center}

\renewcommand{\arraystretch}{1.2}
\begin{longtable}{|L{2.8cm}|L{7.9cm}|L{3.0cm}|}
\hline
\textbf{Component} & \textbf{Heading/Title} & \textbf{Style and Form} \\
\hline
I. Title Pages &
Main title page\newline
Table of contents\newline
List of tables\newline
List of figures\newline
List of abbreviations / key working definitions\newline
List of appendices, if any\newline
Abstract written as one flowing paragraph but covering background, objective, methodology, and expected results or outcome &
Follow the official IUB synopsis format for front matter and the thesis-style formatting rules referenced in the guideline summary chart \\
\hline
II. Main Manuscript &
Introduction\newline
Objectives of the study\newline
Review of literature under suitable subheadings, including a required subheading \emph{Work Conducted in Pakistan}\newline
Justification / rationale of the study\newline
Hypothesis or hypotheses of the study, where applicable\newline
Significance of the study\newline
Novelty statement\newline
Statement of problem\newline
Methodology and plan of work, including geo-location of study, location of sampling, location of laboratory where relevant, study subjects, population or sample number, study attributes or parameters, and data analyses\newline
Likely benefits or outcomes of the study &
Keep the section order close to the IUB pattern and apply the same heading discipline used for formal thesis writing \\
\hline
III. Literature Cited &
Literature cited in APA style &
Use APA formatting consistently and match in-text citations with the final list \\
\hline
\end{longtable}
\renewcommand{\arraystretch}{1.0}

%=============================================================
\section{How to Read the Pattern Correctly}
%=============================================================

\subsection{What the Pattern Is Really Saying}

\begin{itemize}[leftmargin=*]
  \item A synopsis is not only a short proposal; it is a formally organized academic document.
  \item Front matter must be complete and ordered.
  \item The main manuscript must move logically from problem to method to expected contribution.
  \item Literature cited is not optional and must follow APA style.
\end{itemize}

\subsection{What Students Usually Miss}

\begin{itemize}[leftmargin=*]
  \item They write a topic description instead of a problem statement.
  \item They skip \emph{Work Conducted in Pakistan} in the literature review.
  \item They treat methodology as a vague paragraph instead of a plan of work.
  \item They forget novelty, expected benefits, or reference-list consistency.
\end{itemize}

\begin{examplebox}[Minimum Logic of a Good Synopsis]
If a reader asks three questions --- \emph{What is the problem? Why does it matter? How will you study it?} --- the synopsis should answer all three quickly and clearly.
\end{examplebox}

%=============================================================
\section{Guide to the Title Pages}
%=============================================================

\subsection{Main Title Page}

The title page should present the topic, degree program, institutional affiliation, student details, and supervisor details in the official format.

\begin{examplebox}[Title Construction Tip]
A strong IT title is specific, contextual, and researchable. Example: \emph{Adoption of an AI-Based Academic Writing Support System for Synopsis Preparation among MSIT Students at IUB.}
\end{examplebox}

\subsection{Table of Contents, Lists, and Abbreviations}

Include preliminary pages in the order required by the IUB structure. If your synopsis includes tables, figures, abbreviations, or appendices, list them properly.

\subsection{Abstract}

The abstract should be non-structured in appearance but should still touch four ideas:

\begin{itemize}[leftmargin=*]
  \item background of the study,
  \item objective,
  \item methodology,
  \item expected results or outcome.
\end{itemize}

\begin{conceptbox}[Abstract Rule]
Even when the abstract is a single paragraph, it should still behave like a compact summary of the whole study.
\end{conceptbox}

%=============================================================
\section{Guide to the Main Manuscript}
%=============================================================

\subsection{Introduction}

The introduction frames the domain, the institutional or practical context, and the issue that makes the study necessary.

\subsection{Objectives of the Study}

State one general objective and then list specific objectives that are concrete and measurable.

\subsection{Review of Literature}

Organize the review under meaningful subheadings. One subheading must explicitly be:

\begin{center}
\textbf{Work Conducted in Pakistan}
\end{center}

This local-context subheading shows that the researcher has reviewed national or regional work and understands the setting in which the study will be proposed.

\subsection{Justification / Rationale}

Explain why the study is necessary now, who benefits from it, and why existing solutions or existing evidence are inadequate.

\subsection{Hypothesis of the Study}

If the study is quantitative or correlational, state the null and alternative hypotheses clearly. If the study is qualitative, hypotheses may be omitted, but that choice should be methodologically justified.

\subsection{Significance of the Study}

Show the likely academic, practical, institutional, or policy value of the research.

\subsection{Novelty Statement}

Identify what is new. Novelty may come from the population, context, model, dataset, intervention, instrument adaptation, or combination of methods.

\subsection{Statement of Problem}

The statement of problem should be sharper than the introduction. It should identify the precise gap or unresolved condition that the study addresses.

\subsection{Methodology and Plan of Work}

This is not a generic methods paragraph. It is the operational blueprint of the study.

\begin{longtable}{L{3.7cm}L{7.7cm}}
\toprule
\textbf{Subheading} & \textbf{What to Clarify} \\
\midrule
Geo-location of study & The institutional or geographic setting of the study \\
Location of sampling & Where participants or records will be selected \\
Location of laboratory & Where testing, prototyping, or system evaluation will occur if applicable \\
Study subjects & Who or what is being studied \\
Population / sample number & How many participants or records are included \\
Study attributes / parameters & Variables, indicators, or performance measures \\
Data analyses & Statistical or qualitative analysis procedures \\
\bottomrule
\end{longtable}

\subsection{Likely Benefits / Outcomes}

State expected academic, technical, or operational outcomes. Since a synopsis is prepared before the full study is conducted, these should be framed as likely or expected outcomes rather than final findings.

%=============================================================
\section{Guide to Literature Cited in APA Style}
%=============================================================

\begin{definitionbox}[APA Consistency]
APA style in a synopsis means that every in-text citation should have a matching entry in the literature cited section, and every entry in the literature cited section should correspond to a source actually used in the manuscript.
\end{definitionbox}

\begin{examplebox}[Illustrative APA-Style Entries]
The following examples are for formatting demonstration only:

Author, A. A., \& Author, B. B. (2024). Title of the journal article. \emph{Journal Name, 12}(3), 45--61.

Researcher, C. C. (2023). \emph{Title of the book}. Publisher.

Organization Name. (2023). \emph{Title of report}. URL
\end{examplebox}

%=============================================================
\section{Demo Synopsis: Illustrative Sample}
%=============================================================

\begin{conceptbox}[Demo Synopsis Scope]
The following demo synopsis is illustrative. It is written to show sequence, tone, and section depth. Students should replace the sample topic, institutional details, sample size, variables, and references with their own approved study information.
\end{conceptbox}

\subsection{Sample Topic}

\begin{examplebox}[Demo Topic]
\textbf{Adoption of an AI-Based Academic Writing Support System for Synopsis Preparation among MSIT Students at IUB}
\end{examplebox}

%=============================================================
\section{Demo Front Matter}
%=============================================================

\subsection{Demo Main Title Page Layout}

\begin{center}
\begin{tcolorbox}[colback=white,colframe=primary,width=0.88\textwidth,arc=2pt,boxrule=0.8pt]
\begin{center}
{\Large\bfseries Adoption of an AI-Based Academic Writing Support System for Synopsis Preparation among MSIT Students at IUB}\\[12pt]
Synopsis submitted in partial fulfillment of the requirements of the course \emph{Topics in IT}\\[10pt]
Program: MSIT\\
Department / Faculty: [Insert official department name]\\[10pt]
Student Name: [Insert name]\\
Registration No.: [Insert registration number]\\[10pt]
Supervisor: [Insert supervisor name]\\[10pt]
The Islamia University of Bahawalpur\\
2026
\end{center}
\end{tcolorbox}
\end{center}

\subsection{Suggested Preliminary Page Order}

\begin{enumerate}[leftmargin=*]
  \item Main title page
  \item Table of contents
  \item List of tables
  \item List of figures
  \item List of abbreviations / key working definitions
  \item List of appendices, if any
  \item Abstract
\end{enumerate}

\subsection{Demo Abstract}

\begin{examplebox}[Illustrative Abstract]
Graduate IT students increasingly use AI-assisted writing tools during academic preparation, yet their role in synopsis development remains unclear in local university settings. This study aims to examine the factors affecting adoption of an AI-based academic writing support system among MSIT students at IUB and to assess whether perceived usefulness, ease of use, and trust influence intention to use the system for synopsis preparation. The study will employ a quantitative survey design, supported by a short open-ended response section, and data will be collected from MSIT students enrolled during the current academic session. Descriptive statistics and correlation-based analysis are expected to identify the strongest adoption drivers. The likely outcome is a structured understanding of how AI support tools may improve academic planning while also highlighting concerns related to reliability, ethics, and dependency.
\end{examplebox}

%=============================================================
\section{Demo Main Manuscript}
%=============================================================

\subsection{Introduction}

The growing use of artificial intelligence tools in higher education has changed how students search for information, structure arguments, and draft academic content. In graduate IT programs, students are often required to prepare synopses that demand clarity of problem identification, literature review discipline, and methodological planning. However, many students struggle with narrowing the topic, writing the rationale, and organizing the required sections in a formal academic format. AI-based academic writing support systems may help students manage these tasks more efficiently, yet their adoption is shaped by perceptions of usefulness, trust, usability, and academic acceptability. In the context of MSIT students at IUB, the issue is relevant because synopsis quality directly affects research readiness and proposal approval.

\subsection{Objectives of the Study}

\textbf{General Objective:}

To examine the adoption of an AI-based academic writing support system for synopsis preparation among MSIT students at IUB.

\textbf{Specific Objectives:}

\begin{enumerate}[leftmargin=*]
  \item To determine the level of awareness of AI-based writing support tools among MSIT students.
  \item To assess perceived usefulness, ease of use, and trust regarding such tools.
  \item To examine the relationship between these perceptions and intention to use the tools for synopsis preparation.
  \item To identify expected benefits and concerns regarding academic use of AI writing support systems.
\end{enumerate}

\subsection{Review of Literature}

\subsubsection{AI Tools in Academic Writing}

Recent academic discussion suggests that AI-assisted writing environments can support brainstorming, language refinement, structure development, and rapid drafting. At the same time, these tools raise concerns about overreliance, originality, transparency, and accuracy. In synopsis writing, their value may be strongest in outlining sections, organizing ideas, and clarifying academic tone.

\subsubsection{Technology Adoption in Higher Education}

Studies of educational technology adoption commonly emphasize perceived usefulness, perceived ease of use, and trust. These factors influence whether students see a system as worth using for meaningful academic tasks. In the context of synopsis preparation, adoption is likely to depend on whether students believe the tool saves time without weakening academic credibility.

\subsubsection{Work Conducted in Pakistan}

Local higher education research has increasingly focused on digital learning systems, online assessment, and acceptance of educational technologies. Within Pakistani university settings, students often face time constraints, uneven academic writing support, and limited access to structured research mentoring. These conditions make AI-assisted support tools attractive, but local evidence on their role in synopsis preparation remains limited, especially for MSIT-level research preparation.

\subsection{Justification / Rationale of the Study}

The study is justified because synopsis preparation is a critical transition point between coursework and formal research. If AI-based writing support systems are already being used informally, understanding the factors behind their adoption becomes important for academic policy, supervision practices, and student support design. The study is also justified by the need to distinguish helpful academic support from inappropriate dependence.

\subsection{Hypotheses of the Study}

\[H_0^1: \text{Perceived usefulness has no significant relationship with intention to use the system.}\]
\[H_1^1: \text{Perceived usefulness has a significant relationship with intention to use the system.}\]

\[H_0^2: \text{Perceived ease of use has no significant relationship with intention to use the system.}\]
\[H_1^2: \text{Perceived ease of use has a significant relationship with intention to use the system.}\]

\[H_0^3: \text{Trust has no significant relationship with intention to use the system.}\]
\[H_1^3: \text{Trust has a significant relationship with intention to use the system.}\]

\subsection{Significance of the Study}

This study may benefit students by clarifying when and how AI writing support can be used responsibly. It may benefit supervisors by identifying common adoption motivations and concerns. It may also support university-level policy discussions regarding ethical AI use in graduate academic writing.

\subsection{Novelty Statement}

The proposed study is novel because it focuses specifically on AI-based writing support in the narrow and high-stakes context of synopsis preparation among MSIT students. Rather than studying general technology use, it targets a concrete academic task where structure, originality, and methodology are central.

\subsection{Statement of Problem}

Although AI-based writing support tools are becoming more visible in academic environments, their adoption for synopsis preparation among MSIT students at IUB is not yet clearly understood. Students may perceive these tools as useful for organizing ideas and improving academic expression, but concerns about trust, originality, and academic integrity may influence actual adoption. Without evidence on these factors, supervisors and institutions lack a sound basis for guidance and policy.

\subsection{Methodology and Plan of Work}

\textbf{Research Design:}

The study will use a quantitative cross-sectional survey design with a small open-ended response component for contextual clarification.

\textbf{Geo-location of Study:}

The study will be conducted at The Islamia University of Bahawalpur.

\textbf{Location of Sampling:}

Participants will be sampled from MSIT classes and student communication channels approved for academic coordination.

\textbf{Location of Laboratory:}

If system-based demonstration is required, it will be conducted in the departmental computer laboratory or through an approved online demonstration environment.

\textbf{Study Subjects:}

The subjects of the study will be MSIT students engaged in course work or synopsis preparation.

\textbf{Population / Sample Number:}

The population will consist of currently enrolled MSIT students. A sample of approximately 80--120 students may be targeted depending on enrollment and access approval.

\textbf{Study Attributes / Parameters:}

The main variables will include perceived usefulness, perceived ease of use, trust, awareness, and intention to use.

\textbf{Instrument:}

Data will be collected using a structured questionnaire with Likert-scale items and one short open-ended question.

\textbf{Data Analyses:}

The data will be analyzed using frequencies, percentages, mean scores, and correlation-based analysis. Open-ended responses will be summarized thematically.

\textbf{Plan of Work:}

\begin{enumerate}[leftmargin=*]
  \item Finalize the research instrument.
  \item Obtain supervisor and institutional approval.
  \item Collect data from the selected MSIT students.
  \item Clean and code the responses.
  \item Conduct descriptive and relational analysis.
  \item Prepare findings, discussion, and recommendations.
\end{enumerate}

\subsection{Likely Benefits / Outcomes of the Study}

The likely outcome of the study is a clearer understanding of how AI-based academic writing tools are perceived in synopsis development. The results may support student training, supervisor guidance, and policy direction related to ethical AI use in graduate-level academic work.

%=============================================================
\section{Demo Literature Cited (APA Style Illustration)}
%=============================================================

\begin{examplebox}[Illustrative References for Demo Purposes]
The entries below illustrate APA layout only. Replace them with real sources actually used in your synopsis.

Author, A. A. (2023). \emph{Academic writing technologies in higher education}. Example Publisher.

Author, B. B., \& Author, C. C. (2024). Student trust in AI-assisted writing systems. \emph{Journal of Educational Technology Studies, 18}(2), 55--72.

University Research Cell. (2023). \emph{Guidelines for proposal/synopses and theses preparation}. Example University Press.
\end{examplebox}

%=============================================================
\section{Quick Checklist for Students}
%=============================================================

\begin{longtable}{L{9.5cm}C{2cm}}
\toprule
\textbf{Checklist Item} & \textbf{Done?} \\
\midrule
Have I written a precise title instead of a broad topic? & [ ] \\
Have I included all required title pages? & [ ] \\
Does my abstract cover background, objective, method, and expected outcome? & [ ] \\
Have I included \emph{Work Conducted in Pakistan} in the literature review? & [ ] \\
Are my objectives aligned with my problem statement? & [ ] \\
If needed, have I written hypotheses clearly? & [ ] \\
Does my methodology mention setting, subjects, sampling, variables, and analysis? & [ ] \\
Have I explained significance, novelty, and expected benefits? & [ ] \\
Are my in-text citations and literature cited list consistent in APA style? & [ ] \\
\bottomrule
\end{longtable}

%=============================================================
\section{Key Takeaways}
%=============================================================

\begin{itemize}[leftmargin=*]
  \item The synopsis should follow a disciplined academic order, not an informal proposal style.
  \item The IUB pattern requires complete front matter, a structured main manuscript, and APA-style literature cited.
  \item The section \emph{Work Conducted in Pakistan} is a required organizing idea in the literature review under this pattern.
  \item A demo synopsis is useful only if students adapt it honestly to their own approved topic.
\end{itemize}

\vfill
\begin{center}
{\large\textcolor{primary}{\textbf{End of Handout}}}\\[4pt]
{\small\textcolor{gray!70!black}{Synopsis Guide and Demo Synopsis --- Topics in IT, MSIT --- Based on IUB Guidelines 2023}}
\end{center}

\end{document}